\documentclass[11pt, a4paper]{article}
\usepackage[margin=2cm]{geometry}
\usepackage{enumitem}
\usepackage{booktabs}
\usepackage{hyperref}
\usepackage{graphicx}
\usepackage{tabularx}
\usepackage{longtable}
\usepackage{array}
\usepackage{titling}
\usepackage{listings}
\setlength{\parskip}{2pt}
\setlist{nosep, leftmargin=*}

\lstset{
    basicstyle=\ttfamily\footnotesize,
    breaklines=true,
    frame=single,
    language=Python
}

\hypersetup{
    colorlinks=true,
    linkcolor=blue,
    urlcolor=blue
}

\title{\textbf{Assignment 5: Quality Attributes \& Architecture}\\[0.5em]
\large Software Project --- Spring 2026}
\author{}
\date{}

\begin{document}
\maketitle
\thispagestyle{empty}

\begin{center}
\begin{tabular}{ll}
\toprule
\textbf{Project} & PoetryForYou --- Telegram Poetry Learning Bot \\
\textbf{Customer} & Nursultan Askarbekuly \\
\textbf{Team Member} & Egor Zhukov (Solo Developer) \\
\bottomrule
\end{tabular}
\end{center}

% ==================== SPRINT 5 ====================
\section{Sprint 5 Activities}

\subsection*{Sprint Goal}
Improve software quality through architecture documentation, quality attribute scenarios, extended test coverage, and CI pipeline enhancements.

\subsection*{What I Did This Sprint}
\begin{enumerate}
    \item Created issue templates (User Story, Bug Report, Technical Task) and PR template.
    \item Added 5 more unit tests and 5 more integration tests (total: 22 unit, 12 integration).
    \item Added \texttt{mypy} static type checking to the CI pipeline.
    \item Updated README with Development, QA, Build/Deploy, Architecture, and Secrets sections.
    \item Documented quality attribute scenarios using ISO 25010.
    \item Created architecture diagrams (component, sequence, deployment views).
    \item Fixed voice search flow: added ``searching'' stage so voice messages bypass LLM intent classification.
    \item Fixed deployment issues: DB healthcheck, Whisper cache, openai package, model ID.
\end{enumerate}

% ==================== TEMPLATES ====================
\section{Issue \& PR Templates}

Created three issue templates in \texttt{.github/ISSUE\_TEMPLATE/}:

\begin{itemize}
    \item \textbf{User Story} (\texttt{user\_story.md}): Follows ``As a \ldots I want \ldots so that \ldots'' format with acceptance criteria (GIVEN/WHEN/THEN), priority, and story points.
    \item \textbf{Bug Report} (\texttt{bug\_report.md}): Steps to reproduce, expected/actual behavior, environment info.
    \item \textbf{Technical Task} (\texttt{technical\_task.md}): For refactoring/infra work with implementation plan and acceptance criteria.
\end{itemize}

PR template (\texttt{.github/PULL\_REQUEST\_TEMPLATE.md}) includes a checklist: tests, lint, bandit, no secrets, README update, local Docker test.

% ==================== TESTS ====================
\section{Extended Test Suite}

\subsection*{New Unit Tests (5 added)}
\begin{longtable}{p{4cm}p{10cm}}
\toprule
\textbf{Test Class} & \textbf{What It Tests} \\
\midrule
TestChunkSplitting & Splitting poems into learning chunks by stanzas \\
TestFirstLines & Preview text generation for search results \\
TestSearchRanking & Title matches ranked above author-only matches \\
TestLanguageValidation & All language codes produce valid translations \\
TestNormalizeEdgeCases & Cyrillic, mixed-language, and special character handling \\
\bottomrule
\end{longtable}

\subsection*{New Integration Tests (5 added)}
\begin{longtable}{p{4cm}p{10cm}}
\toprule
\textbf{Test Class} & \textbf{What It Tests} \\
\midrule
TestVoiceEndpoint & Voice endpoint validation: missing audio (422), oversized file (413) \\
TestCleanupEndpoint & Session cleanup endpoint returns count \\
TestSearchCommand & \texttt{/search} sets correct stage and returns search prompt \\
TestProfileCommand & \texttt{/profile} returns user profile info \\
TestReviewNoPeems & \texttt{/review} with no learned poems returns helpful message \\
\bottomrule
\end{longtable}

% ==================== CI ====================
\section{CI Pipeline Improvements}

Added \texttt{mypy} as a third static analysis tool (alongside \texttt{ruff} and \texttt{bandit}):

\begin{lstlisting}[language=bash]
# CI pipeline steps:
1. ruff check       # Linting & style
2. bandit scan      # Security vulnerabilities
3. mypy check       # Static type checking (NEW)
4. pytest + coverage # Unit & integration tests
\end{lstlisting}

The mypy step checks \texttt{main.py} and \texttt{schemas.py} with \texttt{--ignore-missing-imports} to catch type errors in the core API layer.

% ==================== README ====================
\section{README Updates}

Added the following sections to \texttt{README.md}:

\begin{itemize}
    \item \textbf{Architecture}: System diagram, technology stack table, component descriptions.
    \item \textbf{Development}: Local setup instructions, project structure tree.
    \item \textbf{Quality Assurance}: CI pipeline table, test coverage summary, local test commands.
    \item \textbf{Build \& Deployment}: Docker commands, VPS deployment guide, auto-restart config.
    \item \textbf{Secrets Management}: Environment variable table, security practices.
\end{itemize}

% ==================== QUALITY ATTRIBUTES ====================
\section{Quality Attribute Scenarios (ISO 25010)}

\subsection*{1. Performance Efficiency --- Time Behavior}

\begin{tabularx}{\textwidth}{lX}
\toprule
\textbf{Source} & Telegram user \\
\textbf{Stimulus} & Sends a text message (e.g., \texttt{/library}) \\
\textbf{Environment} & Production Docker on VPS (1 vCPU, 2 GB RAM) \\
\textbf{Artifact} & Backend API (\texttt{/message} endpoint) \\
\textbf{Response} & Return JSON response \\
\textbf{Measure} & Response time $\leq$ 500ms for text commands \\
\bottomrule
\end{tabularx}

\textit{How achieved}: FastAPI with Uvicorn provides async handling. Database queries use SQLModel with indexes on \texttt{telegram\_id}. All hardcoded poem searches are in-memory ($<$10ms).

\subsection*{2. Security --- Confidentiality}

\begin{tabularx}{\textwidth}{lX}
\toprule
\textbf{Source} & External attacker \\
\textbf{Stimulus} & Attempts to access secrets or database \\
\textbf{Environment} & Production Docker deployment \\
\textbf{Artifact} & Docker containers, source code, `.env` file \\
\textbf{Response} & Secrets remain protected \\
\textbf{Measure} & Zero secrets in source code; container runs as non-root \\
\bottomrule
\end{tabularx}

\textit{How achieved}: Secrets stored in \texttt{.env} (gitignored). Docker containers use \texttt{read\_only: true} filesystem. Backend runs as \texttt{appuser} (non-root). \texttt{bandit} scans in CI detect hardcoded credentials. Database is not exposed to host network.

\subsection*{3. Usability --- Accessibility}

\begin{tabularx}{\textwidth}{lX}
\toprule
\textbf{Source} & User who speaks Russian but cannot type quickly \\
\textbf{Stimulus} & Sends a voice message instead of text \\
\textbf{Environment} & Telegram mobile app \\
\textbf{Artifact} & Voice recognition pipeline (Whisper STT) \\
\textbf{Response} & Transcribe voice and process as text command \\
\textbf{Measure} & Voice recognition accuracy $\geq$ 85\% for Russian speech \\
\bottomrule
\end{tabularx}

\textit{How achieved}: \texttt{faster-whisper} with the \texttt{small} model supports both Russian and English. Voice messages are transcribed server-side and fed into the same message handler. The ``searching'' stage ensures voice input during \texttt{/search} is treated as a search query without LLM re-classification.

\subsection*{4. Reliability --- Availability}

\begin{tabularx}{\textwidth}{lX}
\toprule
\textbf{Source} & System administrator \\
\textbf{Stimulus} & Server reboots unexpectedly \\
\textbf{Environment} & VPS with Docker \\
\textbf{Artifact} & All Docker containers \\
\textbf{Response} & Containers restart automatically \\
\textbf{Measure} & Recovery time $\leq$ 60 seconds after reboot \\
\bottomrule
\end{tabularx}

\textit{How achieved}: Docker Compose uses \texttt{restart: unless-stopped} for all services. PostgreSQL data is persisted in a named volume (\texttt{poetry\_pgdata}). \texttt{systemctl enable docker} ensures Docker starts on boot.

% ==================== ARCHITECTURE ====================
\section{Architecture Views}

\subsection*{Static View: Component Diagram}

\begin{lstlisting}[language={}]
+------------------+         +-------------------+
|   Telegram API   |         |   OpenRouter API  |
+--------+---------+         +--------+----------+
         |                            |
+--------v---------+         +--------v----------+
|    Bot Container  |  HTTP   |  Backend Container|
|  (python-telegram |-------->|   (FastAPI)       |
|   -bot)           |         |                   |
+------------------+         +---+---+---+-------+
                                 |   |   |
                    +------------+   |   +----------+
                    |                |              |
             +------v-----+  +------v------+ +-----v------+
             | PostgreSQL  |  | Whisper STT | | LLM Client |
             |   (DB)      |  | (in-process)| | (OpenRouter)|
             +-------------+  +-------------+ +------------+
\end{lstlisting}

\subsection*{Dynamic View: Voice Search Sequence}

\begin{lstlisting}[language={}]
User -> Telegram: sends voice message
Telegram -> Bot: update with voice file_id
Bot -> Telegram: download voice file
Bot -> Backend: POST /voice (audio + telegram_id)
Backend -> Whisper: transcribe audio
Whisper --> Backend: transcribed text
Backend -> Backend: check user stage ("searching")
Backend -> Backend: extract_search_keywords(text)
Backend -> Backend: search hardcoded + external poems
Backend --> Bot: JSON {reply, suggested_replies}
Bot -> Telegram: send formatted results
Telegram -> User: display poem options
\end{lstlisting}

\subsection*{Deployment View}

\begin{lstlisting}[language={}]
+------------------------------------------+
|            VPS (Ubuntu 24.04)            |
|  +------------------------------------+ |
|  |         Docker Engine               | |
|  |  +----------+ +----------+ +-----+ | |
|  |  |   bot    | | backend  | |  db  | | |
|  |  | :polling | |  :8000   | | :5432| | |
|  |  +----+-----+ +----+-----+ +--+--+ | |
|  |       |             |          |     | |
|  |       +------+------+    +-----+     | |
|  |              |            |          | |
|  |        docker-network   pgdata vol  | |
|  +------------------------------------+ |
+------------------------------------------+
\end{lstlisting}

% ==================== SECRETS ====================
\section{Secrets Management}

All secrets are managed via environment variables in \texttt{.env}:

\begin{longtable}{p{4.5cm}p{1.5cm}p{8cm}}
\toprule
\textbf{Variable} & \textbf{Required} & \textbf{Description} \\
\midrule
TELEGRAM\_BOT\_TOKEN & Yes & Bot token from @BotFather \\
OPENROUTER\_API\_KEY & Yes & API key for LLM features \\
POSTGRES\_USER & Yes & Database username \\
POSTGRES\_PASSWORD & Yes & Database password \\
POSTGRES\_DB & Yes & Database name \\
WHISPER\_MODEL & No & Whisper model size (default: small) \\
LLM\_MODEL & No & LLM model ID \\
\bottomrule
\end{longtable}

\textbf{Security practices}:
\begin{itemize}
    \item \texttt{.env} is in \texttt{.gitignore} --- never committed to Git.
    \item \texttt{.env.example} contains only placeholder values.
    \item Docker containers run as non-root (\texttt{appuser}).
    \item Backend container uses \texttt{read\_only: true} with \texttt{tmpfs} mounts.
    \item \texttt{bandit} scans in CI detect any hardcoded credentials.
\end{itemize}

% ==================== MVP V2 ====================
\section{MVP v2 Deployment}

MVP v2 is deployed on an Aeza VPS and includes:
\begin{itemize}
    \item Voice search with Whisper STT
    \item LLM-powered intent classification and chat
    \item Spaced repetition memorization (SM-2 algorithm)
    \item Library with 25+ classic poems
    \item PostgreSQL persistent storage
\end{itemize}

\textbf{Demo video}: \textit{[TODO: Add link to 2-min demo video]}

% ==================== AI USAGE ====================
\section{AI Usage Report}

\begin{longtable}{p{4cm}p{10cm}}
\toprule
\textbf{Task} & \textbf{AI Contribution} \\
\midrule
Issue/PR templates & Generated templates based on best practices \\
Unit tests & AI generated test structure; manually verified assertions \\
Integration tests & AI generated API test scenarios \\
CI pipeline & AI suggested mypy integration \\
README & AI generated section structure; manually reviewed content \\
Quality scenarios & AI helped format ISO 25010 scenarios; content based on actual architecture \\
Architecture diagrams & AI generated ASCII diagrams from codebase analysis \\
Report & AI generated LaTeX structure; content from actual project state \\
\bottomrule
\end{longtable}

\end{document}
